% ------------------------------------------------------------------------------
% §9 Discussion & policy — drafted Phase C from:
%   findings_task5b.md; findings_task11 (era split, lag-wedge); findings_task13;
%   plan §9; memory: project_sa_regime_exit_2025
% ------------------------------------------------------------------------------
\section{Discussion and Policy Implications}\label{sec:discussion}

\subsection{The conduct is a property of the design}

The mechanism that generates this paper's findings can be stated in one
sentence: South Australia's direction scheme paid gross output at a
backward-looking scarcity price to units whose market alternative was running
at a loss. Each clause contributes. \emph{Gross} payment severs compensation
from the counterfactual, so a unit whose bids establish a low counterfactual
loses nothing by establishing it. The \emph{backward-looking} reference price
held the margin over fuel cost above \$100/MWh for eleven consecutive months
after the 2022 crisis --- peaking at \$211/MWh --- by formula. And the \emph{loss-making market alternative} --- spot
above fuel cost in 3.9\% of directed intervals --- removes any opportunity
cost of absence. Under these payoffs, a standing withdrawn posture dominates
commitment across essentially all conditions, which is exactly the shape of
conduct the mechanism sections document: absence as the resting state,
indifferent to the daily envelope and the daily prize, with depth across
months that tracks what essentiality pays.

The distinction between a standing posture and daily optimisation matters for
how the conduct should be judged. Nothing in this record shows a desk
manipulating the system's need day by day; the informative daily signals ---
event-timed exits, unmoved floor prices, a roster that never falls below the
security standard's ask --- are consistent with a fixed policy plus
housekeeping. The policy itself, however, is priced: the dose response of
Section~\ref{sec:results} is the signature of conduct that sorts on the
payment, and the \$134.7M gap between what the direction channel paid and
what the same megawatt-hours were worth is the transfer that policy earned.
A regulator does not need to prove daily intent to conclude that the scheme
bought its own necessity.

Two scope statements keep the accounting honest. First, the 20-to-1 figure
compares \emph{payment rates on the same megawatt-hours}, not welfare: the
scheme was buying system security, which the spot price of energy does not
value, and the right normative benchmark is the cost of the cheapest
alternative supply of the same service --- the synchronous condensers whose
entry collapsed intervention costs are the revealed measure, and the
counterfactual in which unprofitable units simply retire rather than wait
for directions would have forced that comparison earlier. What the transfer
figure establishes is not a welfare loss of that size but the price
pressure the design placed on a service procurable otherwise. Second, an
owner-level motive this paper cannot fully separate is the gentailer
portfolio: the units' owner retails electricity, and withholding that
raises spot prices pays a portfolio even when it loses the plant money. The
wrong-sign result disciplines this too --- portfolio market power wants
absence in scarce, high-price hours, and the essential intervals where the
absence concentrates are measurably the opposite --- but as a residual
account of the level of withholding (rather than its payment gradient) it
is acknowledged, not excluded.

\subsection{Design alternatives}

Each defect maps to a repair, and the repairs are standard tools of the
capacity-remuneration literature \citep{CramtonOckenfelsStoft2013_capacity,
Fabra2018_capacity}.

\emph{Pay the wedge, not the gross.} Compensating the increment over a
counterfactual dispatch --- or, equivalently, netting retained market revenue
against costs actually incurred --- restores the opportunity-cost margin that
gross payment deletes. The objection that counterfactuals are gameable is
real but misdirected here: the current design already depends on a
counterfactual, the unit's own bid, and pays as if it were zero. The American
reliability-must-run experience documented the same class of defect two
decades earlier: contract forms whose payments were disconnected from market
outcomes created serious incentive problems and raised prices, prompting
redesign \citep{BushnellWolak1999_rmr, BushnellWolak1999_leverage}.

\emph{Look forward, not back.} A trailing-percentile reference price imports
last year's scarcity into this year's payment. Any forward-looking or
cost-anchored reference --- a fuel-indexed price, as this paper's cost-indexed
outcome definition already constructs, or a periodically reset administered
price --- would have cut the 2023 payments roughly in half without touching
the scheme's security function.

\emph{Procure presence explicitly.} The deepest fix prices the actual
commodity. What the system bought was synchronous presence: a known number of
machine-hours per day, from a known set of units. A contracted availability
product converts the rent into a competitive tender and removes the incentive
to be absent at all; the synchronous condensers that ended the bulk direction
era in 2021 (Section~\ref{sec:setting}) are the network-side version of the
same principle --- buy the service, not the emergency.

\subsection{Epilogue: the regime exits}

In September 2025 the minimum synchronous requirement fell to one unit, with
zero targeted once Project EnergyConnect's second stage completes. The
direction channel that paid \$141.9M to three units over 2022--2024 is
closing, not because the conduct was policed but because the underlying
scarcity --- synchronous strength in a renewables-dominated grid --- was
engineered away by transmission and synchronous condensers. That is the final
sense in which the economics of this episode live in the mechanism's design:
the payments were a bridge between an old fleet and a grid that no longer
needs it, and the bridge's toll was set, for three years, by a formula
looking the wrong way.
